\capitulo{7}{Conclusiones y Líneas de trabajo futuras}

En este capítulo final se recogen, por un lado, las principales conclusiones obtenidas tras el desarrollo del proyecto, resaltando los resultados alcanzados y el valor de la solución propuesta, y por otro, algunas posibles líneas de trabajo futuras que permitirían ampliar y reforzar las funcionalidades implementadas.

\section{Conclusiones}

Este Trabajo de Fin de Grado ha permitido desarrollar una aplicación web orientada a la gestión de horarios, reservas y ocupación de aulas en el contexto universitario, integrando en una única plataforma distintas necesidades de planificación académica y administración de espacios. A lo largo del proyecto se han abordado aspectos relevantes tanto de diseño e implementación como de integración de datos, autenticación institucional y gestión de permisos, dando lugar a una solución especializada y adaptada al dominio universitario.

Uno de los elementos más relevantes del desarrollo ha sido la incorporación de un analizador léxico capaz de procesar la información contenida en los horarios académicos y transformarla en reservas periódicas dentro del sistema. Esta funcionalidad, junto con la fase previa de depuración y unificación de la hoja de cálculo original, ha supuesto una de las partes más complejas del proyecto, pero también una de las que más valor aportan a la solución final. En conjunto, el sistema desarrollado constituye una base sólida y funcional para la gestión académica de aulas, superando el enfoque de las herramientas generalistas de reserva de recursos.

\section{Líneas de trabajo futuras}

Como posibles líneas de trabajo futuras, una primera ampliación natural del sistema sería la generación automática de informes del horario y del calendario académico en formatos adecuados para su publicación directa en la web de la universidad. Esto permitiría reforzar la utilidad práctica de la aplicación y mejorar la difusión de la información académica a estudiantes y profesorado.

Asimismo, resultaría interesante incorporar un módulo específico para la gestión de exámenes, de manera que durante los periodos de convocatoria pudiera organizarse también la reserva de aulas asociada a esta actividad. En esta misma línea, podría abordarse el cálculo o propuesta automática de fechas de examen de un curso académico a otro, tomando como base la planificación de años anteriores y facilitando así la preparación inicial del calendario de convocatorias.